\documentclass[11pt,a4paper]{article}

\usepackage[utf8]{inputenc}
\usepackage[T1]{fontenc}
\usepackage[french]{babel}
\usepackage{amsmath,amssymb}
\usepackage{booktabs}
\usepackage{array}
\usepackage{algorithm}
\usepackage{algpseudocode}
\usepackage[margin=2.5cm]{geometry}
\usepackage{hyperref}

\floatname{algorithm}{Algorithme}
\renewcommand{\algorithmicrequire}{\textbf{Entrée :}}
\renewcommand{\algorithmicensure}{\textbf{Sortie :}}

\newcommand{\UB}{\mathrm{UB}}
\newcommand{\CE}{C_E}
\newcommand{\charge}{\mathrm{ch}}

\title{Génération des instances de commandes contrôlées appariées\\(version 2 : arrivées ancrées sur le planning des jobs existants)}
\author{Hedi Boukamcha}
\date{\today}

\begin{document}
\maketitle

\section{Contexte et objectif}

Chaque instance décrit deux commandes. La commande 1 contient les jobs existants $E$, disponibles dès l'instant $0$.
La commande 2 contient les nouveaux jobs $N$, qui arrivent à l'instant de coupure $t_c$ (\emph{cut time}).
À cet instant, la méthode d'acceptation choisit le sous-ensemble de $N$ à insérer dans le planning en cours.

Dans la première version du générateur, $t_c$ était placé par rapport à une borne supérieure $\UB(C^E_{\max})$, qui somme les charges des jobs comme s'ils étaient traités un par un.
Les arrivées des $N$ étaient tirées dans $t_c + \mathcal{U}\{0,\dots,50\}$.
Deux problèmes en découlaient.
\begin{itemize}
  \item La vraie fin des $E$ est souvent bien plus précoce que $\UB(C^E_{\max})$, car les machines et le robot travaillent en parallèle. Dans le tier \emph{late}, environ trois $N$ sur quatre arrivaient après la fin réelle des $E$, dans un atelier déjà vide.
  \item L'étalement de $50$ unités ne dépend pas de la taille de l'instance. Il est large pour les petites instances et négligeable pour les grandes.
\end{itemize}
Les taux d'acceptation mélangeaient donc des insertions réelles et des arrivées après la fin des $E$, ce qui rendait leur analyse incohérente.

La version 2 garantit que \textbf{tous les nouveaux jobs arrivent pendant l'exécution des jobs existants}, et que les tiers \emph{early}, \emph{middle} et \emph{late} ont le même sens quelle que soit la taille.
Pour cela, l'instant de coupure et les arrivées des $N$ sont exprimés en fraction de la fin réelle estimée des $E$, notée $\CE$.
Les opérations, les arrivées des $E$, les dates dues et les coûts sont générés exactement comme dans la version 1.

\section{Notation}

\begin{center}
\begin{tabular}{>{$}l<{$} p{11cm}}
\toprule
\text{Symbole} & Signification \\
\midrule
E,\ N & Jobs existants (commande 1) et nouveaux jobs (commande 2). \\
r_j,\ d_j,\ w_j & Date d'arrivée, date due et coût (poids) du job $j$. \\
p_o & Durée de traitement de l'opération $o$. \\
M = 3,\ L = 2 & Durée d'un déplacement du robot et durée d'un chargement ou déchargement. \\
\charge(j) & Charge d'exécution du job $j$ : $\charge(j) = \sum_{o \in j} (2M + p_o) + 2L$. \\
\UB(C^E_{\max}) & Borne additive de la fin des $E$ (section~\ref{sec:ub}). \\
\UB(C^{E+N}_{\max}) & Borne additive de la fin des $E$ puis des $N$ (section~\ref{sec:ub}). \\
\CE & Fin des $E$ estimée par l'heuristique (section~\ref{sec:heur}). \\
t_c & Instant de coupure, auquel la commande 2 est révélée. \\
\theta & Position relative dans le planning des $E$, avec $t = \theta\,\CE$. \\
\bottomrule
\end{tabular}
\end{center}

\section{Vue d'ensemble}

Pour chaque taille d'instance, le générateur produit des instances d'entraînement et de test.
Chaque instance structurelle $i$ est déclinée selon :
\begin{itemize}
  \item trois ratios $E/N$ : $5/5$, $3/7$ et $7/3$ ;
  \item deux familles de coûts : \texttt{same\_costs} et \texttt{diff\_costs} ;
  \item trois tiers d'arrivée : \emph{early}, \emph{middle} et \emph{late}.
\end{itemize}
Une instance structurelle donne donc $3 \times 2 \times 3 = 18$ fichiers.
L'algorithme~\ref{alg:global} résume la génération.

\begin{algorithm}[ht]
\caption{Génération d'une instance structurelle $i$}
\label{alg:global}
\begin{algorithmic}[1]
\State Tirer $n \sim \mathcal{U}\{n_{\min},\dots,n_{\max}\}$ selon la taille, puis l'ensemble commun de $n$ jobs (section~\ref{sec:pool})
\State Tirer le paramètre $a \sim 10\cdot\mathcal{U}\{0,\dots,10\}$
\For{chaque ratio $\rho \in \{0.5,\ 0.3,\ 0.7\}$}
  \State $E \gets$ les $n_E(\rho)$ premiers jobs ; $N \gets$ les autres \Comment{section~\ref{sec:split}}
  \State Dater les $E$ : arrivées du tirage commun, dates dues par $\UB(C^E_{\max})$ \Comment{section~\ref{sec:due}}
  \State $\CE \gets$ \Call{FinHeuristique}{$E$} \Comment{section~\ref{sec:heur}}
  \For{chaque tier $\tau \in \{\textit{early},\ \textit{middle},\ \textit{late}\}$}
    \State Tirer $t_c$ et la fenêtre d'arrivée $[t_c,\ t_{\mathrm{fin}}]$ \Comment{section~\ref{sec:cut}}
    \State Placer les arrivées des $N$ dans la fenêtre \Comment{section~\ref{sec:arr}}
    \State Dater les $N$ par $\UB(C^{E+N}_{\max})$ \Comment{section~\ref{sec:due}}
  \EndFor
  \For{chaque famille de coûts}
    \State Tirer les coûts et écrire les trois fichiers (un par tier) \Comment{section~\ref{sec:costs}}
  \EndFor
\EndFor
\end{algorithmic}
\end{algorithm}

\section{Étapes détaillées}

\subsection{Ensemble commun de jobs}
\label{sec:pool}

Le nombre de jobs dépend de la taille :
\begin{center}
\begin{tabular}{lcccc}
\toprule
Taille & s & m & l & xl \\
\midrule
Nombre de jobs $n$ & 3 à 5 & 7 à 10 & 15 à 20 & 30 à 50 \\
\bottomrule
\end{tabular}
\end{center}
Chaque job reçoit, comme dans la version 1 :
\begin{itemize}
  \item un nombre d'opérations dans $\mathcal{U}\{1, 2\}$, de types alternés entre la machine 1 et la machine 2 ;
  \item une durée $p_o \sim \mathcal{U}\{10,\dots,60\}$ par opération ;
  \item un indicateur de grande pièce $\mathit{big} \sim \mathcal{U}\{0,1\}$ et un temps de positionnement de $5$ ;
  \item une arrivée relative $r^{\mathrm{rel}}_j \sim \mathcal{U}\{0,\dots,50\}$, utilisée si le job devient un $E$.
\end{itemize}
La version 2 ajoute un tirage $u_j \sim \mathcal{U}(0,1)$ par job.
Il fixe la position d'arrivée du job s'il devient un $N$.
Il est commun aux trois ratios et aux trois tiers, ce qui préserve l'appariement.

\subsection{Découpage entre jobs existants et nouveaux}
\label{sec:split}

Pour un ratio $\rho$, le nombre de jobs existants est
\[
  n_E(\rho) = \min\!\Big(n-1,\ \max\big(1,\ \lfloor n\rho + 0.5 \rfloor\big)\Big).
\]
Les $n_E$ premiers jobs de l'ensemble commun deviennent les $E$, les autres les $N$.
Les $E$ du ratio $3/7$ sont donc inclus dans ceux du ratio $5/5$, eux-mêmes inclus dans ceux du ratio $7/3$.

\subsection{Bornes additives}
\label{sec:ub}

Ces bornes sont reprises de la version 1. Elles ne servent plus qu'aux dates dues.
Les jobs sont parcourus par arrivée croissante :
\begin{align*}
  \UB(C^E_{\max}) &:\quad t \gets 0;\quad \forall j \in E:\ t \gets \max(t, r_j) + \charge(j), \\
  \UB(C^{E+N}_{\max}) &:\quad t \gets \max\big(t_c, \UB(C^E_{\max})\big);\quad \forall j \in N:\ t \gets \max(t, r_j) + \charge(j).
\end{align*}

\subsection{Estimation de la fin réelle des jobs existants}
\label{sec:heur}

$\CE$ est obtenu en ordonnançant les $E$ seuls avec le simulateur utilisé par l'agent, suivi de la recherche locale existante.
L'heuristique est fixe et déterministe. Les instances ne dépendent donc pas de l'agent évalué ensuite.

\paragraph{Représentation d'un ordre.}
Un ordre est une suite de numéros de jobs, où la $k$-ième apparition du job $j$ désigne sa $k$-ième opération.
L'ordre interne de chaque job est ainsi toujours respecté, et les opérations de jobs différents peuvent s'entrelacer.

\paragraph{Évaluation d'un ordre.}
Le simulateur place chaque opération au plus tôt.
Il respecte les arrivées, les déplacements du robot, les chargements et déchargements, les stations, la disponibilité des machines et le positionnement de la machine 1.
La recherche locale teste ensuite, opération par opération, l'activation ou la désactivation du traitement en parallèle de la machine 2 pendant la machine 1.
Elle conserve tout changement qui ne dégrade pas la somme des temps de séjour plus le Cmax.
Pendant cette étape, les dates dues sont fixées égales aux arrivées, pour que le résultat ne dépende d'aucune date due.

\begin{algorithm}[ht]
\caption{\textsc{FinHeuristique}$(E)$}
\label{alg:heur}
\begin{algorithmic}[1]
\Require jobs existants $E$ avec leurs arrivées
\Ensure estimation $\CE$
\State $\mathcal{O} \gets$ \{FIFO par arrivée, SPT, LPT\} \Comment{opérations d'un même job consécutives}
\State Initialiser un générateur aléatoire local avec la graine $1000\,|E| + \sum_{j \in E} r_j$
\For{$k = 1, \dots, 20$}
  \State Ajouter à $\mathcal{O}$ une permutation aléatoire de la suite des opérations
\EndFor
\State $\CE \gets \min_{\sigma \in \mathcal{O}} C_{\max}\big(\textsc{RechercheLocale}(\textsc{Simuler}(\sigma))\big)$
\State \Return $\CE$
\end{algorithmic}
\end{algorithm}

La graine locale rend $\CE$ reproductible et n'altère pas le tirage du reste de l'instance.
Le nombre d'ordres possibles croît très vite avec le nombre d'opérations.
Les $23$ essais couvrent tous les ordres pour les plus petites instances, mais seulement une partie dès la taille m.

\subsection{Instant de coupure et fenêtre d'arrivée}
\label{sec:cut}

Chaque tier est un intervalle $[\theta_{\mathrm{deb}}, \theta_{\mathrm{fin}}]$ du planning des $E$.
Deux paramètres encadrent la fenêtre d'arrivée des $N$ :
\begin{itemize}
  \item une marge $\mu = 0.85$, la fenêtre se termine au plus tard à $\mu\,\CE$ ;
  \item une largeur minimale $\lambda = 0.10$, la fenêtre fait au moins $\lambda\,\CE$.
\end{itemize}
Pour chaque tier :
\begin{align*}
  \theta_{\mathrm{arr}} &= \min(\theta_{\mathrm{fin}},\ \mu), \\
  \theta_c &\sim \mathcal{U}\Big(\theta_{\mathrm{deb}},\ \min\big(\theta_{\mathrm{deb}} + \tfrac{\theta_{\mathrm{fin}} - \theta_{\mathrm{deb}}}{2},\ \theta_{\mathrm{arr}} - \lambda\big)\Big), \\
  t_c &= \max\big(1,\ \mathrm{round}(\theta_c\,\CE)\big), \qquad t_{\mathrm{fin}} = \lfloor \theta_{\mathrm{arr}}\,\CE \rfloor.
\end{align*}
L'instant de coupure est ainsi tiré dans la première moitié de son tier, ce qui évite une fenêtre d'arrivée quasi nulle.
Les valeurs obtenues sont :
\begin{center}
\begin{tabular}{lcccc}
\toprule
Tier & $[\theta_{\mathrm{deb}}, \theta_{\mathrm{fin}}]$ & Intervalle de $\theta_c$ & Fin de la fenêtre $\theta_{\mathrm{arr}}$ & Largeur de la fenêtre \\
\midrule
early  & $[0.05,\ 1/3]$ & $[0.05,\ 0.19]$ & $1/3$  & $0.14$ à $0.28$ \\
middle & $[1/3,\ 2/3]$  & $[0.33,\ 0.50]$ & $2/3$  & $0.17$ à $0.33$ \\
late   & $[2/3,\ 1]$    & $[0.67,\ 0.75]$ & $0.85$ & $0.10$ à $0.18$ \\
\bottomrule
\end{tabular}
\end{center}
Toutes les quantités sont en fraction de $\CE$. Ces valeurs peuvent légèrement varier sur les petites instances, à cause de l'arrondi.

\subsection{Arrivées des nouveaux jobs}
\label{sec:arr}

Chaque nouveau job $j \in N$ arrive à
\[
  r_j = t_c + \mathrm{round}\big(u_j\,(t_{\mathrm{fin}} - t_c)\big),
\]
avec le tirage $u_j$ de la section~\ref{sec:pool}.
Par construction, $r_j \le \mu\,\CE < \CE$ : aucun nouveau job n'arrive après la fin estimée des $E$.
Le générateur vérifie cette propriété pour chaque instance.

La marge $\mu = 0.85$ couvre l'écart entre l'heuristique et l'agent.
Sur les instances de la version 1, l'agent termine les $E$ au plus tôt à $91\,\%$ de l'heuristique en taille s, et à $87\,\%$ en taille m.

\subsection{Dates dues}
\label{sec:due}

Les dates dues gardent la structure de la version 1, avec un plafond plus large et sans garde-fou.
Pour un job d'arrivée $r_j$, on note $\underline{d}_j = \max(b_{\min},\ \max(r_j, 1) + \charge(j))$ la date due la plus précoce admissible, c'est-à-dire la fin du job traité seul.
La date due est tirée par
\[
  d_j \sim \mathcal{U}\Big(\underline{d}_j,\ \max\big(\underline{d}_j,\ b_{\max}\big)\Big),
\]
puis arrondie à l'entier. Les bornes dépendent du type de job :
\begin{center}
\begin{tabular}{lcc}
\toprule
Jobs & $b_{\min}$ & $b_{\max}$ \\
\midrule
Existants $E$ & $50$ & $\UB(C^E_{\max})$ \\
Nouveaux $N$ & $t_c + 50$ & $\UB(C^{E+N}_{\max})$ \\
\bottomrule
\end{tabular}
\end{center}
La date due d'un job se situe donc entre sa fin s'il était traité seul et la borne sérielle de tout le travail.
Si $b_{\max} < \underline{d}_j$, la date due vaut $\underline{d}_j$ et la marge est nulle.
Le facteur multiplicatif de $b_{\max}$ (ici $1$) est réglable par la variable d'environnement \texttt{GNRS\_DUE\_FACTOR}.

\subsection{Coûts}
\label{sec:costs}

\begin{center}
\begin{tabular}{lcc}
\toprule
Famille & Coût des $E$ & Coût des $N$ \\
\midrule
\texttt{same\_costs} & $1$ & $1$ \\
\texttt{diff\_costs} & $\mathcal{U}(0.8,\ 1.0)$ & $\mathcal{U}(0.05,\ 0.4)$ \\
\bottomrule
\end{tabular}
\end{center}
Les coûts sont arrondis à quatre décimales.
Le coût d'un nouveau job est identique dans les trois tiers.

\section{Appariement}

L'appariement permet de comparer les scénarios instance par instance.
\begin{itemize}
  \item \texttt{same\_costs} et \texttt{diff\_costs} partagent exactement les mêmes jobs, arrivées, dates dues et instants de coupure. Seuls les coûts diffèrent.
  \item Les trois ratios partagent le même ensemble de jobs. Seul le découpage entre $E$ et $N$ change, ainsi que tout ce qui en dépend : $\CE$, $t_c$, arrivées des $N$ et dates dues.
  \item Les trois tiers partagent les mêmes $E$, les mêmes coûts et les mêmes tirages $u_j$. Seuls $t_c$, les arrivées des $N$ et leurs dates dues changent.
\end{itemize}

\section{Changements par rapport à l'ancien générateur}
\label{sec:changes}

L'ancien générateur est \texttt{instance\_generator\_paired.py}. Ses instances sont dans \texttt{data/controlled\_orders\_paired}.
Le nouveau générateur est \texttt{instance\_generator\_paired\_v2.py}. Ses instances sont dans \texttt{data/controlled\_orders\_paired\_v2}.
L'ancien générateur et ses instances n'ont pas été modifiés.

\subsection{Ce qui change}

\paragraph{1. La référence temporelle.}
\begin{itemize}
  \item \emph{Avant :} les tiers étaient définis sur $\UB(C^E_{\max})$, qui additionne les charges des $E$ comme s'ils passaient un par un.
  \item \emph{Maintenant :} les tiers sont définis sur $\CE$, la fin des $E$ obtenue en les ordonnançant réellement avec le simulateur et la recherche locale (section~\ref{sec:heur}).
  \item \emph{Pourquoi :} les machines et le robot travaillent en parallèle. La vraie fin des $E$ est donc souvent de $5$ à $35\,\%$ plus précoce que $\UB(C^E_{\max})$ en taille m, et l'écart grandit avec la taille. De plus, $\UB(C^E_{\max})$ n'est pas toujours une vraie borne supérieure, car elle ignore le temps de positionnement de la machine 1. Un instant de coupure placé sur cette borne ne correspondait donc pas à l'avancement réel des $E$.
\end{itemize}

\paragraph{2. L'instant de coupure.}
\begin{itemize}
  \item \emph{Avant :} $t_c$ était tiré uniformément dans tout le tier, soit $[0, \tfrac13]$, $[\tfrac13, \tfrac23]$ ou $[\tfrac23, 1]$ fois $\UB(C^E_{\max})$.
  \item \emph{Maintenant :} $t_c$ est tiré dans la première moitié du tier, en fraction de $\CE$, en laissant une fenêtre d'arrivée d'au moins $0.10\,\CE$ (section~\ref{sec:cut}). Le tier \emph{early} commence à $0.05\,\CE$.
  \item \emph{Pourquoi :} un instant de coupure tiré près de la fin du tier laissait une fenêtre d'arrivée presque nulle. La position de $t_c$ devenait alors une source de variabilité non contrôlée.
\end{itemize}

\paragraph{3. Les arrivées des nouveaux jobs.}
\begin{itemize}
  \item \emph{Avant :} $r_j = t_c + \mathcal{U}\{0,\dots,50\}$, quelle que soit la taille de l'instance.
  \item \emph{Maintenant :} $r_j$ est uniforme entre $t_c$ et $\min(\theta_{\mathrm{fin}}, 0.85)\,\CE$ (section~\ref{sec:arr}). La position relative $u_j$ de chaque job est commune aux ratios et aux tiers.
  \item \emph{Pourquoi :} avec l'ancien tirage, environ $75\,\%$ des $N$ du tier \emph{late} et $15\,\%$ de ceux du tier \emph{middle} arrivaient après la fin réelle des $E$. Ces jobs trouvaient un atelier vide et étaient presque toujours acceptés. Les taux d'acceptation mélangeaient ainsi de vraies insertions et des arrivées sans concurrence. L'étalement fixe de $50$ unités rendait aussi les tailles incomparables : il était large pour les petites instances et négligeable pour les grandes.
\end{itemize}

\paragraph{4. Les dates dues.}
\begin{itemize}
  \item \emph{Avant :} $d_j \sim \mathcal{U}\big(\underline{d}_j,\ \max(\underline{d}_j + 50,\ \UB/3)\big)$, avec $\UB = \UB(C^E_{\max})$ pour les $E$ et $\UB(C^{E+N}_{\max})$ pour les $N$.
  \item \emph{Maintenant :} $d_j \sim \mathcal{U}\big(\underline{d}_j,\ \max(\underline{d}_j,\ \UB)\big)$. Le plafond passe de $\UB/3$ à $\UB$, et le garde-fou $\underline{d}_j + 50$ est supprimé.
  \item \emph{Pourquoi :} le plafond $\UB/3$ tombait le plus souvent sous $\underline{d}_j$. En taille s, c'était le cas de $97\,\%$ des $N$ ; en taille m, de $68$ à $96\,\%$ selon le tier. La date due était alors fixée par le garde-fou, une marge arbitraire entre $0$ et $50$, et non par la formule. Supprimer le garde-fou sans élargir le plafond donnait une marge nulle à $91\,\%$ des $E$ et à $96$ à $99\,\%$ des $N$ \emph{middle} et \emph{late} en taille s. Avec le plafond $\UB$, la formule fixe réellement la date due : la part de marges nulles tombe sous $2\,\%$ partout, sauf pour les $E$ de taille s ($14\,\%$, cas d'un seul job existant, sans concurrence entre $E$).
\end{itemize}

\paragraph{5. Les métadonnées des instances.}
Chaque fichier JSON contient trois nouveaux champs : \texttt{generator}, \texttt{cmax\_e\_heuristic} et \texttt{arrival\_window}.
Le reste du format est identique, et les instances se lisent avec le même code.

\subsection{Ce qui ne change pas}

\begin{itemize}
  \item Les tailles, le nombre d'instances par taille et la graine par défaut.
  \item Les opérations des jobs : nombre, types alternés, durées, indicateur de grande pièce et temps de positionnement.
  \item Les arrivées des $E$, tirées dans $\mathcal{U}\{0,\dots,50\}$.
  \item Le découpage entre $E$ et $N$ selon le ratio.
  \item La forme des dates dues : tirage uniforme entre la fin au plus tôt du job et une borne additive, avec $b_{\min} = 50$ pour les $E$ et $t_c + 50$ pour les $N$.
  \item Les deux familles de coûts et leurs intervalles.
  \item L'appariement entre familles de coûts, entre ratios et entre tiers.
\end{itemize}

\subsection{Conséquences pratiques}

\begin{itemize}
  \item Les agents entraînés sur l'ancien jeu doivent être réentraînés sur le nouveau, avec l'option \texttt{--dataset=controlled\_orders\_paired\_v2} du script d'entraînement.
  \item Les résultats obtenus sur l'ancien jeu ne sont pas directement comparables à ceux du nouveau. Les instances de même numéro ont les mêmes opérations, mais pas les mêmes instants de coupure, arrivées et dates dues des $N$.
  \item Le tier \emph{late} représente désormais une vraie insertion tardive pendant la production, et non plus une arrivée majoritairement après la fin des $E$.
\end{itemize}

\subsection{Récapitulatif}

\begin{center}
\begin{tabular}{p{3.6cm} p{5.2cm} p{5.2cm}}
\toprule
Élément & Version 1 & Version 2 \\
\midrule
Référence temporelle & $\UB(C^E_{\max})$, borne additive & $\CE$, fin heuristique des $E$ \\
Instant de coupure & Uniforme dans tout le tier de $\UB(C^E_{\max})$ & Première moitié du tier de $\CE$ \\
Largeur de la fenêtre d'arrivée & $50$ unités, fixe & au moins $0.10\,\CE$ \\
Arrivées des $N$ & $t_c + \mathcal{U}\{0,\dots,50\}$ & Uniformes entre $t_c$ et $\min(\theta_{\mathrm{fin}}, 0.85)\,\CE$ \\
$N$ arrivés après la fin des $E$ & environ $75\,\%$ en \emph{late}, $15\,\%$ en \emph{middle} & $0\,\%$ par construction \\
Arrivées des $E$ & $\mathcal{U}\{0,\dots,50\}$ & inchangé \\
Plafond des dates dues & $\max(\underline{d}_j + 50,\ \UB/3)$ & $\max(\underline{d}_j,\ \UB)$, sans garde-fou \\
Coûts et appariement & & inchangé \\
\bottomrule
\end{tabular}
\end{center}

\section{Jeu de données produit et validation}

Le jeu contient $30$ instances d'entraînement et $20$ de test par taille.
Avec les $18$ déclinaisons, cela donne $3600$ fichiers, avec la graine $1$.
\begin{itemize}
  \item Aucun nouveau job n'arrive après $\CE$, sur l'ensemble des $3600$ fichiers.
  \item Les paires \texttt{same\_costs} et \texttt{diff\_costs} sont identiques hors coûts.
  \item Aucune structure de test n'apparaît dans l'entraînement.
\end{itemize}

Chaque fichier JSON garde le format de la version 1, lisible par le solveur et par l'entraînement.
Il ajoute trois champs : \texttt{generator} (\texttt{paired\_v2}), \texttt{cmax\_e\_heuristic} pour $\CE$, et \texttt{arrival\_window} pour $[t_c, t_{\mathrm{fin}}]$.

\section{Limites connues}
\label{sec:limits}

\paragraph{Jobs existants souvent en retard au départ.}
Même avec le plafond $\UB(C^E_{\max})$, au moins un $E$ est souvent déjà en retard dans le planning heuristique de la section~\ref{sec:heur}, avant l'arrivée des $N$. Nombre d'instances concernées sur $20$ instances de test :
\begin{center}
\begin{tabular}{lccc}
\toprule
Taille & $3/7$ & $5/5$ & $7/3$ \\
\midrule
s  & 6  & 13 & 16 \\
m  & 14 & 19 & 20 \\
l  & 19 & 20 & 20 \\
xl & 20 & 20 & 20 \\
\bottomrule
\end{tabular}
\end{center}
C'est structurel : avec un tirage uniforme à partir de $\underline{d}_j$, dès que les $E$ sont nombreux, l'un d'eux reçoit presque sûrement une marge faible et devient en retard au moindre conflit.
Ces valeurs surestiment un peu la réalité, car l'heuristique minimise le $C_{\max}$ et non le retard.
La règle d'acceptation dépend de ce retard initial.
S'il est nul, chaque $E$ peut être retardé d'au plus $\varepsilon\,(d_j - r_j)$.
Sinon, le retard pondéré des $E$ ne doit pas dépasser $(1+\delta)$ fois le retard initial.
Les comparaisons d'acceptation entre ratios doivent donc être faites à régime d'acceptation égal.

\paragraph{Marges nulles résiduelles.}
Sans garde-fou, un job reçoit une marge nulle quand $b_{\max} < \underline{d}_j$. Part des jobs concernés dans le jeu produit :
\begin{center}
\begin{tabular}{lcccc}
\toprule
Taille & $E$ & $N$ early & $N$ middle & $N$ late \\
\midrule
s  & 14.1\,\% & 0\,\% & 0\,\%   & 1.9\,\% \\
m  & 0.7\,\%  & 0\,\% & 0.4\,\% & 0.4\,\% \\
l  & 0\,\%    & 0\,\% & 0\,\%   & 0\,\% \\
xl & 0.1\,\%  & 0\,\% & 0\,\%   & 0.1\,\% \\
\bottomrule
\end{tabular}
\end{center}
Les $14\,\%$ de $E$ en taille s correspondent aux instances à un seul job existant, pour lequel $\UB(C^E_{\max}) = \underline{d}_j$.

\paragraph{Écart entre l'heuristique et l'agent.}
La marge $\mu = 0.85$ couvre les écarts observés en tailles s et m.
Elle n'a pas été vérifiée en tailles l et xl, où les $23$ ordres testés ne couvrent qu'une infime partie des ordres possibles.
Après les tests, la propriété peut être contrôlée en comparant la fin des $E$ dans le planning de l'agent aux arrivées des $N$.

\paragraph{Granularité du taux d'acceptation.}
En taille s, une instance contient de $1$ à $3$ nouveaux jobs, et le taux d'acceptation varie par paliers.
Il est préférable de rapporter des moyennes par scénario avec un intervalle de confiance, ou le nombre total de jobs acceptés, plutôt que des minimums et maximums.

\section{Reproduction}

\begin{verbatim}
python3 instance_generator_paired_v2.py --train=30 --test=20 --path=./
\end{verbatim}
Les instances sont écrites dans
\texttt{data/controlled\_orders\_paired\_v2/<train|test>/<taille>/<scénario>/instance\_<i>\_<tier>.json}.
L'option \texttt{--sizes} restreint la génération à certaines tailles, et \texttt{--seed} change la graine.

\end{document}
